% !TEX root = ./main.tex
% =====================================================
% PREÂMBULO — pacotes, configurações e ambientes
% (complementa o main.tex; NÃO redefine fontes, idioma,
%  margens, cabeçalho/rodapé ou hyperref — isso já está
%  configurado no main.tex)
% =====================================================



% ---------- Matemática ----------
\usepackage{amsmath}
\usepackage{amssymb}
\usepackage{amscd}
\usepackage{amsthm}
\usepackage{mathtools}
\usepackage{bm}               % negrito em símbolos matemáticos

% ---------- Figuras e desenhos ----------
\usepackage{tikz}
\usetikzlibrary{arrows.meta, calc, positioning, decorations.markings, angles, quotes}
\usetikzlibrary{shapes.geometric}
\usetikzlibrary{shapes.symbols}
\usetikzlibrary{shadows}
\usepackage{float}
\usepackage{subcaption}

% ---------- Tabelas ----------
\usepackage{multirow}

% ---------- Referências cruzadas e citações ----------
\usepackage[
  backend=biber,
  style=abnt,
  sorting=nyt
]{biblatex}
\addbibresource{referencias.bib}

\usepackage[capitalise,brazilian]{cleveref}

% =====================================================
% AMBIENTES MATEMÁTICOS
% Numerados por capítulo (ex.: Teorema 2.1, Lema 2.2, ...)
% Todos compartilham o mesmo contador para manter a ordem
% cronológica de aparição no texto.
% =====================================================
\theoremstyle{plain}
\newtheorem{theorem}{Teorema}
\newtheorem{proposition}[theorem]{Proposição}
\newtheorem{lemma}[theorem]{Lema}
\newtheorem{corollary}[theorem]{Corolário}

\theoremstyle{definition}
\newtheorem{example}[theorem]{Exemplo}
\newtheorem{examples}[theorem]{Exemplos}
\newtheorem{definition}[theorem]{Definição}
\newtheorem{axiom}{Axioma}

\theoremstyle{remark}
\newtheorem{observation}[theorem]{Observação}

% Fim de demonstração customizado (símbolo QED)
\renewcommand{\qedsymbol}{$\blacksquare$}

% Nome em minúsculo (uso em \cref) e maiúsculo (uso em \Cref)
\crefname{axiom}{axioma}{axiomas}
\Crefname{axiom}{Axioma}{Axiomas}
\crefname{theorem}{teorema}{teoremas}
\Crefname{theorem}{Teorema}{Teoremas}
\crefname{proposition}{proposição}{proposições}
\Crefname{proposition}{Proposição}{Proposições}
\crefname{lemma}{lema}{lemas}
\Crefname{lemma}{Lema}{Lemas}
\crefname{corollary}{corolário}{corolários}
\Crefname{corollary}{Corolário}{Corolários}
\crefname{example}{exemplo}{exemplos}
\Crefname{example}{Exemplo}{Exemplos}
\crefname{examples}{exemplos}{exemplos}
\Crefname{examples}{Exemplos}{Exemplos}
\crefname{definition}{definição}{definições}
\Crefname{definition}{Definição}{Definições}
\crefname{observation}{observação}{observações}
\Crefname{observation}{Observação}{Observações}